% "{'chapitre':'slci_multiphy','classe':('PSI'),'type':('application'),'titre':'Mise à l\\'eau d\\'un robot sous-marin','source':'Concours Centrale - MP 2019','comp':['B2-02',],'corrige':True}"

%\setchapterimage{fig_00.jpg}
\chapter*{Application \arabic{cptApplication} \\ 
Actionneur ElectroMécanique (EMA) de la tuyère de la fusée VEGA -- \ifprof Corrigé \else Sujet \fi}
\addcontentsline{toc}{section}{Application \arabic{cptApplication} : 
Actionneur ElectroMécanique (EMA) de la tuyère de la fusée VEGA -- \ifprof Corrigé \else Sujet \fi}

\iflivret \stepcounter{cptApplication} \else
\ifprof  \stepcounter{cptApplication} \else \fi
\fi

\setcounter{question}{0}
%\marginnote{Concours Banque PT.}





\subsection*{Question de cours}
\question{Donner la strucutre de la chaîne de puissance. On précisera en particulier le nom de chacune des fonctions.}

\question{Donner 3 exemples de composants permettant de réaliser chacune des fonctions.}

\subsection*{Actionneur ElectroMécanique (EMA) de la tuyère de la fusée VEGA}
\marginnote{D'après ressources de Florestan Mathurin.}


\begin{marginfigure}
\includegraphics[width=\linewidth]{fig_01}
\end{marginfigure}

Le développement du lanceur européen VEGA a démarré en 1998 et s’est achevé en 2011. Ce projet répondait
à une demande de mise en orbite basse et polaire, à coûts réduits, de satellites scientifiques ou d’observation
de la terre, dont la masse peut aller jusqu’à 2000 kg. On s’intéresse ici à l’Actionneur ElectroMécanique (EMA)
de la tuyère de cette fusée.
Lors du programme de développement, la minimisation des coûts
s’est appuyée sur l’implémentation de technologies avancées déjà
disponibles et l’utilisation des installations des lanceurs Ariane.
Cependant certaines parties ont fait l’objet d’innovation comme le
système de contrôle vectoriel de poussée (en Anglais : « Thrust
Vector Control ») du premier étage de propulsion P80. D’une
longueur de dix mètres, le P80 est chargé de 88 tonnes de
propergol solide. Ceci lui permet de disposer d’une poussée
maximale de 3000 kN et d’un temps de combustion de 107
secondes. Afin de bien contrôler la trajectoire de la fusée il est
indispensable d’orienter très rapidement et très précisément la
tuyère du P80. Alors que sur Ariane 5 le pilotage vectoriel de la
poussée est assuré par des dispositifs à source de puissance
hydraulique, sur le P80 cette tâche est assurée par des dispositifs à
source de puissance électrique (en Anglais : « Power By Wire »).
On donne la structure du système de l’actionneur sous la forme
d’un diagramme de définition de blocs (bdd).

\begin{center}
\includegraphics[width=\linewidth]{fig_02}
\end{center}


On donne sur le document réponse DR1 le diagramme de bloc interne (ibd).


\question{Sur ce diagramme de bloc interne DR1, colorier en rouge les flux correspondants à l’énergie électrique de puissance, en bleu les flux correspondants à l’énergie mécanique de puissance et en vert les flux
correspondant à des informations. Il peut rester des traits en noirs !}

\question{Compléter sur le document réponse DR2, le diagramme chaîne de puissance/chaîne d’information
correspondant au système tuyère orientable VEGA.}

\begin{center}
\includegraphics[width=\linewidth]{fig_03}
\end{center}


\begin{center}
\includegraphics[width=\linewidth]{fig_04}
\end{center}


